\chapter{Team Commitment Contract}
\label{tab:TCC}
\section*{\textcolor{Delft_Blue}{\textbf{Team name: B2-1\vspace{5pt}}}}
\section*{\textcolor{Delft_Blue}{\textbf{Shared team values:\vspace{5pt}}}}

\begin{enumerate}
    \item \textbf{Be present (Physically \& mentally).}\\
- \textit{We discussed the importance of being physically present during our team meeting. We all agreed that being physically there will increase our performance, as well as allow us to communicate better. Therefore we put it in the contract, to increase our likelihood of success.} \\
    \item \textbf{Keep your promises.}\\
- \textit{During the team meeting we agreed that by making promises to each other regarding the work, deadlines or anything regarding the project, we will make sure to keep them. This is important to us as it will show that each of us is committed to the project, and it will keep any false promises from happening.}\\
    \item \textbf{Be direct but respectful.}\\
- \textit{We think that the ability to be direct with one another is an important thing, as it will create an environment where open communication can happen. This will be important during the project, as we can then give one another useful tips, or feedback. Additionally, we agreed that being respectful is another key factor, as it will keep open communication from being hurtful towards others.}\\
    \item \textbf{Integrity.}\\
- \textit{We all agreed that having integrity is a very important thing in this project. Because we are continuously working together, we agreed that we should all keep the same moral standard in our work, so no copying or using artificial intelligence as this will harm the academic integrity of our work.}\\
\end{enumerate}


\section*{\textcolor{Delft_Blue}{\textbf{Target or ambition level:\vspace{5pt}}}}
\textbf{What grade are you working for?}\\
    - \textit{We are all aiming for a 10.0. This is a high grade, but we agreed that setting our goals and expectations high will motivate us to work harder and reach this goal. Even if we do not all reach this grade, striving for a 10.0 will allow us to work to our maximum potential.}\\
\newpage
\section*{\textcolor{Delft_Blue}{\textbf{Quality:\vspace{5pt}}}}
\textbf{How will you prevent fraud and plagiarism?}\\
    - \textit{By regularly asking each other to check our work. By doing this we will not only make sure that the quality of the work is of a high standard, it also ensures that the group mates can spot and point out any potential plagiarism. Moreover, we instilled a policy where we will not copy anything from others or the internet, such as AI. This can therefore help us to prevent plagiarism }\\
    
\textbf{How you will organize feedback (relation, process, content)?}\\
    - \textit{We have two team meetings a week, one on Tuesday and one on Thursday. We will dedicate some time during every team meeting on Thursday for feedback. We will allow people to give their feedback on their sub-group's work and the whole group's work.}\\

\textbf{How do you determine the quality of the work your group and each team member does?}\\
    - \textit{During the project, other group members will read over everyone's work, allowing them to spot and point out any mistakes. Moreover, we will check sources as there are a lot of unreliable sources that can harm the quality of our work. Therefore having quality sources can help to raise the quality of our work. Lastly, we will run simulations when possible, we have tools such as 'LT Spice' or 'Python' available to us, which can help us determine wether the work is up to the standards we have set, using the knowledge we have about the subject.}\\
    

\section*{\textcolor{Delft_Blue}{\textbf{Workload:\vspace{5pt}}}}
\textbf{How do you make sure that the workload is fairly and evenly distributed?\\ }
    - \textit{We decided in a team meeting that this is up to the sub-groups to decide, as the sub-groups are most aware of their workload. This therefore puts the responsibility of an equitable workload on the people in their respective sub-groups. Additionally, we decided that if other group members are aware of the workload not being distributed fairly in a specific sub-group, then this can be called out and action can be taken. This may include a discussion in the next team meeting, and distributing the workload more fairly.}\\

\section*{\textcolor{Delft_Blue}{\textbf{Planning:\vspace{5pt}}}}
\textbf{How do you ensure that each team member finishes everything on time? Did you clarify who will have a final say in the final deliverable and submit it to Brightspace on behalf of the project group?\\}
    - \textit{We also decided that the responsibility of finishing things on time first lies within the sub-groups, however, if deadlines are not being met then the other group members will start to get involved, as they are directly affected by the slacking behavior of other group members. We have decided on a group member who will deliver the final project on Brightspace.}\\
    
\section*{\textcolor{Delft_Blue}{\textbf{Behavior:\vspace{5pt}}}}
\textbf{How do you treat each other in the group? \\}
    - \textit{The most important thing is to treat each other as equals and with respect. By treating each other as equals, we can introduce new ways of thinking into the group and stimulate a safe working environment. By always being respectful we can ensure that we avoid conflicts and maintain a good place for working.}\\

\textbf{How do you handle disagreements/conflicts within your group? \\}
    - \textit{Disagreements are bound to happen, as we are always working on the same project together, and therefore it is very likely that there will be differing opinions. We want to handle these conflicts in an orderly fashion, first allowing the disagreeing people to talk about why and find some sort of middle ground. If this is not possible there will be a third member involved, who can then help find a middle ground.}\\

\textbf{Could your guide or TA be involved in reaching consent?\\}
    - \textit{If there is a conflict, and there has already been a third member involved without being able to find a middle ground, we could ask a TA or guide to help find a middle ground. For inspiration we may also be able to ask TA's for help, however, we do not want the TA's to introduce ideas that we would not have been able to come up with ourselves, as that might be considered inequitable.}\\

\textbf{What do you do if someone is late during a group meeting?\\}
    - \textit{We will first encourage them to explain why they are late, as they might have a valid reason that caused them to be late. If this reason is not valid we will give them a warning. If they continue to be late we will sit down as a group and discuss it in detail. We will explain the harm it does to the team to encourage the person not to be late again. If this does keep happening, the late comer could be assigned any task this person wasn't present for when the group was discussing who is  responsible.}\\


\section*{\textcolor{Delft_Blue}{\textbf{Communication:\vspace{5pt}}}}
\textbf{In what ways do you communicate with each other as a group and among yourselves? Zoom, MS Teams?\\}
- \textit{We will communicate through WhatsApp as this is where we are all most active, which will allow us to answer questions that others have as quickly and as often as possible. Moreover, we will be using Notion for the planning of the project, it has a built-in calendar and a task divider, which will make it clear what each person has to do both daily and weekly.}\\

\textbf{What information do you share via WhatsApp, e-mail, telephone?\\}
    - \textit{We will share most of our information through WhatsApp, as this is the quickest form of communication for all of us. This might include calculations, text for in any group document or Python code. Otherwise, we might use E-Mail for files and Notion for task division and planning.}\\
    
\textbf{What are your office hours?\\} 
    - \textit{We decided that we want the weekdays, Monday to Friday, from 08:30 to 18:00 to be the 'Office Hours'. This is when we should expect questions to be answered quickly, and answer questions ourselves quickly too. During the 'Office Hours' we should all be as active as possible. After 'Office Hours' we will not be expected to answer questions, however, we still can, as we will likely work on documents for the project at the weekends. Team members should understand and respect that someone may not respond after the 'Office Hours' as they have other things to do as well.}\\

\textbf{How do you become an expert on the topic you are not directly involved in?\\}
    - \textit{During the team meetings on Thursday, when we do an update on each of the sub-groups, we expect them to give a detailed description of anything they are working on and explain how it works. This allows the other group members to understand how other components work and also become experts.}\\

\section*{\textcolor{Delft_Blue}{\textbf{Decision-making:\vspace{5pt}}}}
\textbf{How do you make decisions? By majority vote or by consensus?\\}
    - \textit{We will make our decisions by consensus most of the time, as this will please most of the group. It will not be likely that there are too many differing opinions as we all agree on most things. Therefore consensus can give us the best results.}\\

\section*{\textcolor{Delft_Blue}{\textbf{Consequences:\vspace{5pt}}}}
\textbf{What are the consequences if a participant in the group does not keep the agreements?\\}
    - \textit{We will create a strike agreement, if someone gets three strikes they will be expected to bring cake the next day and share it with our group. This is a fun way of making sure that people will be keeping their agreements. Because this approach may not seem too serious we will also make sure that at three strikes, the person who has the strikes will get a serious talking to from the other group members. They will explain the consequences of not keeping the agreements.}\\

\section*{\textcolor{Delft_Blue}{\textbf{Success factors:\vspace{5pt}}}}
\textbf{What makes your team a dream team? \\}

- \textit{What makes us a dream team is that everyone in the group is motivated to work hard and is eager to learn. This creates an optimal environment for group work and allows us to become a good team and excel in the project.}\\
\newpage
\section*{\textcolor{Delft_Blue}{\textbf{Norms or evaluation criteria\vspace{5pt}}}}
You will evaluate your own and each other's work in this project.\\
Discuss and write down your team criteria. You need at least five different criteria. \\

\begin{enumerate}
    \item Keeps deadlines\\
    \item Writes everything themselves\\
    \item Being able to explain what they worked on in detail\\
    \item How well they participated throughout the project\\
    \item The quality of the work\\
\end{enumerate}

\raggedleft
\textbf{Signature of involved team members and date:\\}
\textit{Niek F. van Vroonhoven | 24 - 11 - 2024}
\newline\textit{Alex Guo | 24 - 11 - 2024}
\newline\textit{Stijn J.H. ter Huurne | 25 - 11 - 2024}
\newline\textit{Karsten W. Smid | 25 - 11 - 2024}
\newline\textit{Noah T. Herijgers | 25 - 11 - 2024}
\newline\textit{G.D. van Weelden | 25 - 11 - 2024}
\newline\textit{Itai Givony | 27 - 11 - 2024}
\newline\textit{Teun J. Bello | 28 - 11 - 2024}
\newline\textit{Tycho V. Wever | 28 - 11 - 2024}
\newline\textit{Luuk D. Melissant | 01 - 12 - 2024}

